\chapter{Discussion}
\label{ch:discussion}

This chapter interprets our experimental findings, framing them as systems engineering insights for building latency-sensitive search components on managed runtimes.

\section{GC-Neutral Memory Engineering}
Dynamic runtimes are historically considered unsuitable for performance-critical systems because automatic garbage collection (GC) sweeps introduce unpredictable latency spikes. Kronos shows that these pauses are avoidable by eliminating dynamic object allocations on the hot path.

As documented in Chapter~\ref{ch:experimental_results}, early configurations utilizing uncapped transposition tables suffered from heap inflation, which triggered frequent collection cycles. Restricting transposition table growth and wrapping board mutations in the allocation-free \texttt{ChessInstanceClone} class bounded V8 heap memory during standard tournament matches. Within this footprint, V8's minor collections run in under 1\,ms, preventing search loop interruptions.

This memory discipline enables a sustained search throughput across search depths. Although this throughput is lower than native C++ engines, it is sufficient to evaluate depth-6 horizons within standard time controls. For systems developers working on managed runtimes, the core takeaway is that memory allocation patterns, rather than language choice, govern real-time search efficiency.

\section{Synergy between Move Ordering and Search Reductions}
Our ablation data indicates a multiplicative synergy between ordering heuristics and tree reductions. The cumulative search stack compresses the search space compared to baseline alpha-beta search without ordering and search enhancements. LMR builds upon this sorting to reduce the evaluation depth of late-ranked moves. Disabling LMR increases the search space significantly even with ordering active.

This interaction is highly non-linear. The safety of LMR is conditioned on move ordering quality; if the best move is ranked late in the list, LMR searches it at reduced depth, potentially overlooking a tactical variation. Conversely, poor ordering under an active LMR configuration is more hazardous than poor ordering in a standard alpha-beta search. Thus, investment in move ordering quality (MVV-LVA, Killer Moves, and History Heuristics) directly enhances the safety of aggressive pruning heuristics.

\section{Heuristic Marginal Utility and Search Failures}
Analyzing individual heuristics reveals a stark divide between search-space reduction and playing-strength outcomes:
\begin{itemize}
    \item \textbf{LMR Dominance vs. History Heuristics Marginal Utility:} Late Move Reduction (LMR) acts as the primary driver of node savings in Kronos, compressing the search tree geometrically by searching late-ranked moves at a reduced depth. However, its effectiveness is highly dependent on move sorting quality. In contrast, the History Heuristic exhibits very marginal utility. In a shallow search horizon (depth 3 to 4), the search tree is too small for quiet move beta cutoffs to accumulate statistically representative history scores. Consequently, the history table contains noisy weights that fail to improve move ordering, indicating that history-based sorting requires deeper search horizons to provide positive utility.
    \item \textbf{Null Move Pruning (NMP) Limitations and Pruning Failures:} At shallow horizons (such as depth 4), NMP fails to provide meaningful node reductions. Under our configuration ($R = 2$), a null-move search at depth 4 is executed at depth $d - 1 - R = 1$ ply. A 1-ply search with a null move incurs significant move verification and evaluation overhead, yet rarely triggers beta cutoffs because static positional evaluations at such shallow depths do not fluctuate enough to justify the pruning assumption. This shallow-depth pruning failure is well-documented in computer chess literature~\cite{donninger1993null}: NMP with a constant reduction factor $R = 2$ is widely understood to be practically ineffective below depth 5. Consequently, the final engine configuration activates NMP only for depths $d \ge 5$.
    \item \textbf{SPRT Inconclusiveness and Test Resolution Bounds:} Under the configured SPRT protocol ($H_0: \text{Elo} \le 0$ vs. $H_1: \text{Elo} \ge 10$), only quiescence search crossed the Wald boundary to be accepted. The remaining enhancements (PVS, LMR, NMP, and Killer/History) remained inconclusive. This indicates that while these heuristics succeed in reducing node count and executing searches faster, their individual playing strength contributions are below the statistical resolution of the test ($< 10$ Elo) under the configured sample size. To resolve such minor Elo differentials, a significantly larger sample size (exceeding thousands of games) would be required.
    \item \textbf{Single-Threaded Latency Constraints and V8 JIT Warming:} The performance regression observed at depth 7 is a consequence of single-threaded latency constraints combined with V8 execution mechanics. The V8 engine relies on hotness thresholds (number of function invocations) to trigger JIT compilation (from the Ignition interpreter up to Sparkplug and eventually TurboFan). In the initial stages of a search, code runs interpreted or baseline-compiled. As the search deepens, parallel JIT compilation threads optimize the hot paths, but the synchronization overhead and concurrent compilation tasks introduce transient latency spikes (JIT compilation stutter). At depth 7, these latency spikes cause the synchronous worker thread to exceed the tournament time control limit. The search must abort and fall back to the results of the last fully completed iterative deepening ply (depth 6). This hard latency ceiling indicates that single-threaded JavaScript cannot reliably execute deep game-tree searches under standard competitive time controls.
\end{itemize}

\section{Latency Constraints and the Search Ceiling}
Depth-scaling metrics exhibit consistent strength gains up to depth 6, followed by regression at depth 7. This is a latency-induced limitation of the single-threaded execution model. Additionally, the non-monotonic NPS behavior remains an observed telemetry behavior requiring future investigation; we hypothesize this stems from transposition table cache collision dynamics and JIT recompilation triggers under V8.

At depth 7, the average search latency per position increases substantially, which frequently exceeds the tournament time limit. When the search is aborted mid-ply, the engine reverts to the best move from the last completed iterative deepening step (depth 6). In positions where depth-6 and depth-7 evaluations differ, the engine plays sub-optimally. This indicates a hard throughput ceiling, verifying that single-threaded JavaScript workers cannot reliably sustain depth-7 search under standard time controls. Pushing past this limit would require shared-memory parallel search strategies, such as Lazy SMP, operating over browser \texttt{SharedArrayBuffer} allocations.

\section{Deterministic Seeding and SPRT Validation}
To ensure reproducibility in playing strength evaluation, Kronos incorporates deterministic seeding and sequential testing.

First, the tournament runner shuffles the opening library using a custom LCG PRNG initialized with a fixed seed. This ensures that independent benchmark runs execute identical game sequences. Second, paired, color-flipped matches isolate playing strength from opening advantages. Each opening FEN is played twice, swapping color assignments between variants, which mitigates White's first-move advantage.

Third, Sequential Probability Ratio Testing (SPRT) provides a mathematically grounded termination criterion. Rather than evaluating a fixed game count, SPRT tracks the LLR against Wald boundaries, terminating as soon as the outcomes statistically confirm or reject a rating hypothesis. In practice, this sequential approach reduces the required game sample size by up to 40\% compared to fixed-sample testing, making laptop-based systems evaluations viable.

\textbf{Determinism and statistical validity:} Because the search core is strictly deterministic (no temperature-scaled randomness or stochastic move selection), two identical engine variants playing from the same opening FEN will produce the exact same move sequence every time. The effective number of independent game outcomes per matchup is therefore bounded by the opening book size: $2 \times |\text{opening book}| = 100$ independent observations (50 positions $\times$ 2 color assignments). In cross-variant matchups (e.g., Alpha-Beta vs.\ Minimax), the two engine variants produce different move sequences from the same opening, so each of the 100 starting configurations yields a unique game trajectory. However, once all 100 configurations have been played, subsequent games are exact replays of earlier outcomes. Consequently, the reported SPRT game counts in Table~\ref{tab:statistical_validation} (ranging from 153 to 400 games) exceed this independence bound, meaning the effective sample sizes are capped at 100 regardless of the total games executed. The reported LLR values and confidence intervals are therefore computed over a sample containing duplicated outcomes, which inflates the apparent statistical power. We acknowledge this as a structural limitation of deterministic self-play: the SPRT termination decisions should be interpreted with the understanding that the true effective sample size is at most 100, and the reported confidence intervals understate the actual uncertainty. Crucially, these higher game counts (up to 400 games) were intentionally sustained to evaluate the long-running memory stability and garbage collection resilience of the tournament scheduler and worker threads over hundreds of consecutive games, rather than to claim additional statistical degrees of freedom.

To assess the impact of this sample-size inflation on our empirical findings, as a conservative sensitivity check, we linearly scaled down the reported LLR metrics to the $N_{\text{eff}} = 100$ limit (effectively dividing the LLR by the inflation factor $N_{\text{total}}/100$). For Quiescence Search (EXP-6), scaling the LLR from 400 down to 100 games yields an adjusted LLR of $75.433 / 4 \approx 18.86$. Because this adjusted value remains significantly above the upper Wald boundary of 2.944, the acceptance of the playing-strength hypothesis for quiescence search remains statistically valid. For all other configurations (e.g., Move Ordering, EXP-2, where the LLR scales from 0.859 to $\approx 0.46$; Iterative Deepening, EXP-5, where the LLR scales from $-$0.171 to $\approx -0.11$), the adjusted metrics remain well within the Wald boundaries, confirming their inconclusive status. Therefore, while sample-size correction doubles the width of our rating confidence intervals, it does not alter the final binary acceptance or rejection decisions of the SPRT framework.

\section{Draw-Rate Analysis in Self-Play Pools}
During self-play matchups between variants of similar strength, draw rates reached 66.7\% to 70.0\% (Table~\ref{tab:tournament_results}). We identify four contributing factors to this high draw rate:
\begin{enumerate}
    \item \textbf{Shallow Search Depth:} At depths 3 and 4, the engines cannot formulate long-term positional plans to create structural imbalances, resulting in defensive simplifications.
    \item \textbf{Search Determinism:} Root move selection is entirely deterministic without randomization. In a closed self-play pool, this lacks variation, leading games into identical drawn lines. Combined with the bounded opening book, this determinism means that high draw rates between similarly-ranked variants are structurally inevitable rather than indicative of true playing strength parity.
    \item \textbf{Move Limit Termination:} Even matchups frequently simplify into drawish endgames that are terminated by the runner's 200-ply limit.
    \item \textbf{Repetition Detection Limitations:} The transposition table does not penalize move repetitions during search traversal. Consequently, the search core is prone to repeating moves when it evaluates the position as equal.
\end{enumerate}
This high draw rate reduces the statistical information generated per game, expanding standard error margins and requiring larger game sample sizes to resolve minor Elo differences. In practice, this directly impacted the execution time of the automated tournament scheduler: because draws provide less LLR divergence than decisive results, matchups between variants of similar strength (such as EXP-3 and EXP-4) were forced to run close to or up to their maximum game limits before the LLR could cross a Wald boundary or stabilize within confidence limits.

\section{Managed Runtime Engineering Lessons}
Three design guidelines emerge from our study:
\begin{enumerate}
    \item \textbf{Prioritize Allocation Elimination:} Replacing object allocations with in-place board state mutations yields a greater practical speedup than algorithmic modifications by removing GC spikes.
    \item \textbf{Warm the JIT Compiler:} V8 JIT optimizations require several warm-up iterations to stabilize. Benchmarking timings recorded during the first few matches will result in distorted, high-latency metrics.
    \item \textbf{Isolate Concurrent Tasks:} Without background Web Worker isolation, search execution blocks the browser main thread, freezing the user interface. Decoupling the search plane is a functional requirement for interactive dynamic runtimes.
\end{enumerate}
